% ============================================================
%  Chapter 2 — Component Architecture (Internal)
% ============================================================
\chapter{Component Architecture}
\label{ch:architecture}

This chapter describes the internal architecture of \textbf{resilient-call}.
The content here is intended for maintainers and contributors, not end-users.

\section{Design Principles}
\label{sec:design-principles}

The architecture was designed with the following goals in mind:

\begin{itemize}[leftmargin=*, label=\textbullet]
    \item \textbf{Composability over inheritance.} Retry and circuit-breaker
          logic are independent objects that \texttt{ResilientExecutor} composes
          rather than combining in a single monolithic class.  Either mechanism
          can be used alone.
    \item \textbf{Zero dependencies.} The entire library relies only on the
          Python standard library, making installation and auditing
          straightforward.
    \item \textbf{Testability.} All state is explicit and accessible through
          public properties, enabling deterministic unit testing without
          mocking \texttt{time}.
    \item \textbf{Thread safety.} Every mutable state transition is protected
          by a \texttt{threading.Lock}.
    \item \textbf{Extensibility through callbacks.} Rather than subclassing,
          callers extend behaviour via \texttt{on\_retry} and
          \texttt{on\_state\_change} hooks.
\end{itemize}

\section{Module Structure}
\label{sec:module-structure}

\begin{figure}[H]
    \centering
    \dirtree{%
        .1 resilient\_call/.
        .2 \_\_init\_\_.py \DTcomment{public re-exports}.
        .2 exceptions.py \DTcomment{CircuitBreakerError, RetryExhaustedError}.
        .2 circuit\_breaker.py \DTcomment{CircuitBreaker, Config, State enum}.
        .2 retry.py \DTcomment{RetryPolicy, BackoffStrategy enum}.
        .2 executor.py \DTcomment{ResilientExecutor}.
        .2 registry.py \DTcomment{CircuitBreakerRegistry + module singleton}.
        .2 decorators.py \DTcomment{@retry, @circuit\_breaker}.
    }
    \caption{Package module layout}
    \label{fig:module-structure}
\end{figure}

\section{Class Diagram}
\label{sec:class-diagram}

Figure~\ref{fig:class-internal} shows the full internal class diagram including
private members.

\begin{figure}[H]
    \centering
    \includegraphics[width=\linewidth]{class_diagram_internal}
    \caption{Internal class diagram — all members}
    \label{fig:class-internal}
\end{figure}

\subsection{Class Responsibilities}

\begin{table}[H]
    \centering
    \caption{Summary of class responsibilities}
    \label{tab:class-responsibilities}
    \begin{tabularx}{\linewidth}{>{\ttfamily}l X}
        \toprule
        \textbf{Class} & \textbf{Responsibility} \\
        \midrule
        CircuitBreakerState  & Enum of three states: CLOSED, OPEN, HALF\_OPEN. \\
        CircuitBreakerConfig & Value object holding all tuning parameters. \\
        CircuitBreaker       & Thread-safe state machine; wraps a callable and tracks failure/success counts. \\
        BackoffStrategy      & Enum of delay growth algorithms: FIXED, EXPONENTIAL, JITTER. \\
        RetryPolicy          & Executes a callable up to \texttt{max\_attempts} times with computed delays between attempts. \\
        ResilientExecutor    & Composes RetryPolicy and CircuitBreaker; applies fallback on final failure. \\
        CircuitBreakerRegistry & Named store of CircuitBreaker instances; provides a module-level singleton. \\
        CircuitBreakerError  & Raised when a call is rejected because the circuit is OPEN. \\
        RetryExhaustedError  & Raised when all retry attempts fail. \\
        \bottomrule
    \end{tabularx}
\end{table}

\section{State Diagram}
\label{sec:state-diagram}

The \texttt{CircuitBreaker} implements a finite state machine with three states.
Figure~\ref{fig:state-diagram} shows all transitions and their trigger conditions.

\begin{figure}[H]
    \centering
    \includegraphics[width=0.82\linewidth]{state_diagram}
    \caption{Circuit Breaker state machine}
    \label{fig:state-diagram}
\end{figure}

\begin{table}[H]
    \centering
    \caption{State transition conditions}
    \label{tab:transitions}
    \begin{tabular}{llp{7cm}}
        \toprule
        \textbf{From} & \textbf{To} & \textbf{Condition} \\
        \midrule
        CLOSED    & OPEN      & Consecutive failures $\geq$ \texttt{failure\_threshold} \\
        OPEN      & HALF\_OPEN & \texttt{timeout} seconds elapsed since opening \\
        HALF\_OPEN & CLOSED    & Consecutive successes $\geq$ \texttt{success\_threshold} \\
        HALF\_OPEN & OPEN      & Any single failure (timer resets) \\
        \bottomrule
    \end{tabular}
\end{table}

\section{Usage Scenarios — Sequence Diagrams}
\label{sec:sequence-diagrams}

\subsection{Scenario 1: Successful Call}

The service responds on the first attempt.  The circuit breaker records a success
and resets the failure counter (Figure~\ref{fig:seq-success}).

\begin{figure}[H]
    \centering
    \includegraphics[width=0.9\linewidth]{seq_success}
    \caption{Sequence diagram — successful call, no retry needed}
    \label{fig:seq-success}
\end{figure}

\subsection{Scenario 2: Transient Failure, Retries Succeed}

The service fails twice with a transient error and recovers on the third attempt.
Exponential backoff is applied between attempts
(Figure~\ref{fig:seq-retry}).

\begin{figure}[H]
    \centering
    \includegraphics[width=0.9\linewidth]{seq_retry_success}
    \caption{Sequence diagram — two failures followed by successful retry}
    \label{fig:seq-retry}
\end{figure}

\subsection{Scenario 3: Circuit Open, Fallback Returned}

The circuit breaker has already tripped to OPEN.  Calls are rejected immediately
without contacting the downstream service; the configured fallback value is
returned to the caller (Figure~\ref{fig:seq-open}).

\begin{figure}[H]
    \centering
    \includegraphics[width=0.9\linewidth]{seq_circuit_open}
    \caption{Sequence diagram — circuit OPEN, service bypassed, fallback returned}
    \label{fig:seq-open}
\end{figure}

\section{Deployment Diagram}
\label{sec:deployment-diagram}

\textbf{resilient-call} is a pure library component — it has no server-side
counterpart and requires no infrastructure beyond a Python 3.10+ runtime
(Figure~\ref{fig:deployment}).

\begin{figure}[H]
    \centering
    \includegraphics[width=\linewidth]{deployment_diagram}
    \caption{Deployment diagram — library installed via pip, used inside application server}
    \label{fig:deployment}
\end{figure}

Key deployment characteristics:

\begin{itemize}[leftmargin=*, label=\textbullet]
    \item Distributed as a wheel via PyPI (\texttt{pip install resilient-call}).
    \item All state is in-process; no external services are required.
    \item In horizontally scaled deployments each replica holds its own
          circuit breaker state independently.
\end{itemize}

\section{Execution Flow}
\label{sec:execution-flow}

When \texttt{ResilientExecutor.execute(func)} is called the following
order of operations applies:

\begin{enumerate}[leftmargin=*, label=\arabic*.]
    \item \textbf{Circuit check.} The \texttt{CircuitBreaker} checks its
          current state under the lock.  If OPEN, \texttt{CircuitBreakerError}
          is raised immediately.  If the timeout has elapsed, the state
          transitions to HALF\_OPEN.
    \item \textbf{Retry loop.} The \texttt{RetryPolicy} calls \texttt{func}
          up to \texttt{max\_attempts} times, sleeping between attempts
          according to the configured \texttt{BackoffStrategy}.
    \item \textbf{Success recording.} On a successful return, the
          \texttt{CircuitBreaker} records a success (closing the circuit if
          in HALF\_OPEN).
    \item \textbf{Failure recording.} On final failure, the
          \texttt{CircuitBreaker} records a failure (potentially opening
          the circuit), and \texttt{RetryExhaustedError} propagates upward.
    \item \textbf{Fallback.} If a fallback callable is configured on the
          executor, any propagating exception is caught and the fallback
          return value is returned to the caller instead.
\end{enumerate}
